To utilize this approach we first identify the two types of fibrations that will lend themselves to this method. The first is the path space fibration $\Om B\lra E\lra B$ introduced in \cref{ex:pathspacefibr}. Since $E$ is contractible, we apply the long exact sequence in homotopy to see that if $B=K(G,n)$, then $F\simeq K(G,n-1)$.

Our other type of fibration will come from a short exact sequence of groups $A \lra B\lra C$ and will result in a fibration of the form 
\[
K(A,n)\lra K(B,n)\lra K(C,n)
\]
for every $n$. To do this, we will need to define the notion of a homotopy fiber and familiarize ourselves with the fact that every map $f:X\lra Y$ is homotopically equivalent to a Hurewicz. This is formalized in the next proposition. 

\begin{prop}
    Let $f:X\lra Y$ be a map between topological spaces. Then $f$ is homotopically equivalent to a map $f':X'\lra Y$ which is a Hurewicz fibration such that $X'$ deformation retracts onto $X$. 
\end{prop}
\begin{proof}
    
\end{proof}
In light of this proposition, if we have a short exact sequence of abelian groups $0\lra A\lra B\lra C\lra 0$, then we may realize this as a fibration of Eilenberg-MacLane spaces as follows: due to the correspondence between group homomorphisms $B\lra C$ and $[K(B,n)\lra K(C,n)]$, we obtain a map $K(B,n)\lra K(C,n)$ which we may convert into a fibration $F\lra K(B,n)\lra K(C,n)$ and apply the long exact sequence in homotopy. We have
\[
\pi_{n+1}(K(C,n))\lra \pi_n(F)\lra \pi_n(K(B,n))\lra \pi_n(K(C,n))\lra \pi_{n-1}(F)\lra\cdots
\]
\textbf{explain how to conclude this bit}

We will now compute some fundamental examples with these fibrations, some of which are already known, but nonetheless are important to illustrate how this spectal sequence works operationally. 

\begin{ex}
    $H^*(\Z,2;\Z)=H^*(\CP^{\infty})$, fibration is $S^1=K(\Z,1)\lra E \lra K(Z,2)$
\end{ex}

\begin{ex}
    $H^*(\Z/2,1;\Z/2)$ using the fibration coming from the short exact sequence $0\lra \Z\lra \Z\lra\Z/2\lra 0$
\end{ex}

We will now move on to calculate the cohomology of a space that one might not see in an introductory algebraic topology class. We compute $H^*(\Z/2,2;\Z/2)$ in low dimensions using the path space fibration $K(\Z/2,1)\lra * \lra K(\Z/2,2)$. We know the cohomology of the total space, since it is contractible, and of the fiber by the previous example, which is $\Z/2[\a]$ where $\a$ is a 1 dimensional class. Coefficients in $\Z/2$ are assumed and will be suppressed from notation.

The spectral sequence for this fibration with $\Z/2$ coefficients converges to $H^*(*)$, so we see that $E^{p,q}_{\infty}=0$ except at $(0,0)$ where it is equal to $\Z/2$. We will use the following diagram to aid us in our computations:

\textbf{diagram}

Since $B$ is simply connected, the Hurewicz theorem tells us that $H^1(B)=0$, so that $E_r^{1,q}=0$ for all $q$ and $r$. Since in particular this means that $d_1(\a)=0$ we see that $\a$ is transgressive and that $\tau(\a)=d_2(\a)=\iota_2$, which generates a $\Z/2$ at $E^{2,0}_2$, meaning $H^2(B)=\Z/2$. Noting that $H^2(F)=\Z/2$ generated by $\a^2$, we use the fact that $d_2$ is a derivation to note that $d_2(\a^2)=d_2(\a)\otimes \a - \a \otimes d_2(\a) =0$. \textbf{check} In fact, we can notice that $d_2(\a^{2k})=0$ for any $k$ by an induction argument. On the other hand, $d_2(\a^{2k+1})=d_2(\a)\otimes \a^{2k}=\iota_2\otimes \a^{2k}$. Since $E_2$ is a tensor product of $H^*(B)$ and $H^*(F)$, we can see that 